\documentclass[12 pt, letterpaper]{hmcpset}
\usepackage{hw-preamble}

\name{Jayden Thadani}
\class{MATH 139 --- Fourier Analysis}
\assignment{HW \#1}
\duedate{29th January 2025}

\begin{document}

\begin{problem}[2.]
	If $z = x + i y$ is a complex number with $x, y \in \mathbb{R}$, we define the \emph{complex conjugate} of $z$ by
	\[
		\overline{z} = z - i y.
	\]
	\begin{enumerate}
		\item What is the geometric interpretation of $\overline{z}$.
		\item Show that $\lvert z \rvert^{2} = z \overline{z}$.
		\item Prove that if $z$ belongs to the unit circle, then $1 / z = \overline{z}$.
	\end{enumerate}
\end{problem}

\begin{solution}
	\begin{enumerate}
		\item Conjugation corresponds to reflection in the real axis.
		\item We can just compute both sides for $z = x + i y$.
			\begin{align*}
				z \overline{z} & = (x + i y) (x - i y) \\
					       & = x^{2} - i^{2} y^{2} + i x y - i x y \\
					       & = x^{2} + y^{2} \\
					       & = \lvert z \rvert^{2}.
			\end{align*}
		\item If $z$ belongs to the unit circle, $\lvert z \rvert = 1$. Note that (for $z \neq 0$)
			\[
				z \frac{\overline{z}}{\lvert z \rvert^{2}} = \frac{\lvert z \rvert^{2}}{\lvert z \rvert^{2}} = 1.
			\]
			Since inverses are unique, we have  $1 / z = \overline{z} / \lvert z \rvert^{2}$. For $z$ on the circle where $\lvert z \rvert^{2} = 1$, this is just saying $1 / z = \overline{z}$.
	\end{enumerate}
\end{solution}

\pagebreak

\begin{problem}[3.]
	\begin{enumerate}
		\item Show that a converging sequence of complex numbers has a unique limit.
		\item Prove that a sequence of complex numbers converges if and only if it is a Cauchy sequence.
		\item  Let $\{ a_{n} \}_{n = 1}^{\infty}$ be a sequence of non-negative real numbers such that the series $\sum_{n} a_{n}$ converges. Show that if $\{ z_{n} \}_{n = 1}^{\infty}$ is a sequence of complex numbers satisfying $\lvert z_{n} \rvert \le a_{n}$ for all $n$, then the series $\sum_{n} z_{n}$ converges.
	\end{enumerate}
\end{problem}

\begin{enumerate}
	\item Suppose there is a convergent sequence $\{ z_{n} \}_{n \in \mathbb{N}}$ that converges to both, $z$ and $z'$. We will show that $z = z'$. 

		By definition, for any real $\varepsilon > 0$, there are $N, N' \in \mathbb{N}$ such that for all $n > N$ we have $\lvert z - z_{n} \rvert < \varepsilon / 2$ and for $n > N'$ we have $\lvert z' - z_{n} \rvert < \varepsilon / 2$. Thus, there is some $N'' \in \mathbb{N}$ such thatfor all $n > N''$, we have both $\lvert z - z_{n} \rvert < \varepsilon / 2$ and $\lvert z' - z_{n} \rvert < \varepsilon / 2$.

		By the triangle inequality, this means that for any $\varepsilon > 0$ we have
		\begin{align*}
			\lvert z - z' \rvert & \le \lvert z - z_{n} \rvert + \lvert z_{n} - z' \rvert \\
					     & < \varepsilon/2 + \varepsilon / 2 \\ 
					     & = \varepsilon.
		\end{align*}
		Since $z$ and $z'$ have distance less than any $\varepsilon > 0$, they must have distance $0$, and thus, be the same number --- $z = z'$.
	\item Suppose that $\{ z_{n} \}_{n \in \mathbb{N}}$ is a sequence of complex numbers converging to $z$. Again, for any real $\varepsilon > 0$, there is some $N$ such that for each pair $n, m > N$ we have $\lvert z_{m} - z \rvert < \varepsilon / 2$ and $\lvert z_{n} - z \rvert < \varepsilon / 2$. By the triangle inequality, we get
		\begin{align*}
			\lvert z_{m} - z_{n} \rvert & \le \lvert z_{m} - z \rvert + \lvert z - z_{n} \rvert \\ 
						    & < \varepsilon / 2 + \varepsilon / 2 \\
						    & = \varepsilon.
		\end{align*}
		That is, for any $\varepsilon > 0$, there is some $N \in \mathbb{N}$ such that for each pair $m, n > N$ we have $\lvert z_{m} - z_{n} \rvert < \epsilon$ --- the sequence $\{ z_{n} \}_{n \in \mathbb{N}}$ is Cauchy.

		Suppose $\{ z_{n} = x_{n} + i y_{n} \}_{n \in \mathbb{N}}$ is a Cauchy sequence of complex numbers. We will show that the sequences of real numbers $\{ x_{n} \}_{n \in \mathbb{N}}$ and $\{ y_{n} \}_{n \in \mathbb{N}}$ are Cauchy, and further, for $x_{n} \to x$ and $y_{n} \to y$, we have $z_{n} \to z = x + i y$.

		Note that for any $\varepsilon > 0$, there is some $N \in \mathbb{N}$ such that for  each pair $m, n > N$ we have $\lvert z_{n} - z_{m} \rvert < \varepsilon$. Equivalently,
		\[
			(x_{n} - x_{m})^{2} + (y_{n} - y_{m})^{2} < \varepsilon^{2}.
		\]
		Thus, for each $\varepsilon > 0$ there exists some $N \in \mathbb{N}$, such that for each pair $m, n > N$ we have both $\lvert x_{n} - x_{m} \rvert < \varepsilon$ and $\lvert y_{n} - y_{m} \rvert < \varepsilon$. That is, both $\{ x_{n} \}_{n \in \mathbb{N}}$ and $\{ y_{n} \}_{n \in \mathbb{N}}$ are Cauchy sequences.

		Any Cauchy sequence of real numbers converges ($\mathbb{R}$ is a complete metric space). Say $x_{n} \to x$ and $y_{n} \to y$ as $n \to \infty$. Then consider $z = x + i y$. For each real $\varepsilon > 0$, there exists some $N \in \mathbb{N}$ such that for all $n > N$, both $\lvert x_{n} - x \rvert$ and $\lvert y_{n} - y \rvert$ are less than $\varepsilon / \sqrt{2}$. Thus, for all $n > N$ we have
		\begin{align*}
			\lvert z_{n} - z \rvert^{2} & = (x_{n} - x)^{2} + (y_{n} - y)^{2} \\
						    & < \varepsilon^{2}/2 + \varepsilon^{2} /2 \\
						    & = \varepsilon^{2}.
		\end{align*}
		Thus, for all $\varepsilon > 0$ there is some $N \in \mathbb{N}$ such that for all $n > N$, we have $\lvert z_{n} - z \rvert < \varepsilon$ --- $z_{n} \to z$  as $n \to \infty$.
	\item If the series $\sum_{n} a_{n}$ converges, then it's partial sums are a convergent, and thus, Cauchy sequence. That is, for each $\varepsilon > 0$ there is some $N \in \mathbb{N}$ such that each pair $m, n$ gives $\sum_{k = m + 1}^{n} a_{k} < \varepsilon$.

		Since $\lvert z_{n} \rvert \le a_{n}$, the series $\sum_{n} z_{n}$ has partial sums satisfying
		\begin{align*}
			\left \lvert \sum_{k = 1}^{n} z_{k} - \sum_{k = 1}^{m} z_{k} \right \rvert & = \left \lvert \sum_{k = m + 1}^{n} z_{n} \right \rvert \\ 
												   & \le \sum_{k = m + 1}^{n} \lvert z_{n} \rvert \\
												   & = \sum_{k = m + 1}^{n} a_{n} \\
												   & < \varepsilon.
		\end{align*}
		That is, the partial sums of the series are Cauchy, and thus, the series converges.
\end{enumerate}

\pagebreak

\begin{problem}[4.]
	For $z \in \mathbb{C}$ we define the \emph{complex exponential} by
	\[
		e^{z} = \sum_{n = 0}^{\infty} \frac{z^{n}}{n!}
	\]
	\begin{enumerate}
		\item Prove that the above definition makes sense, by showing that the series converges for every complex number $z$. Moreover, show that the convergence is uniform on every bounded subset of $\mathbb{C}$.
		\item If $z_{1}, z_2$ are two complex numbers, prove that $e^{z_{1}} e^{z_{2}} = e^{z_{1} + z_{2}}$.
		\item Show that if $z$ is purely imaginary, that is $z = i y$ with $y \in \mathbb{R}$, then
			\[
				e^{z} = e^{i y} = \cos{y} + i \sin{y}.
			\]
			This is Euler's identity.
		\item More generally,
			\[
				e^{x + iy} = e^{x} (\cos{y} + i \sin{y})
			\]
			whenever $x, y \in \mathbb{R}$, and show that
			\[
				\lvert e^{x + i y} \rvert = e^{x}.
			\]
		\item Prove that $e^{z} = 1$ if and only if $z = 2 \pi k i$ for some integer $k$.
		\item Show that every complex number $z = x + i y$ can be written in the form
			\[
				z = r e^{i \theta}
			\]
			where $r$ is unique and in the range $0 \le r < \infty$, and $\theta \in \mathbb{R}$ is unique up to an integer multiple of $2 \pi$. Check that
			 \[
				r = \lvert z \rvert \quad \text{and} \quad \theta = \arctan{(y / x)}
			\]
			whenever these formulas make sense.
		\item In particular $i = e^{i \pi / 2}$. What is the geometric meaning of multiplying a complex number by $i$? Or by $e^{i \theta}$ for any $\theta \in \mathbb{R}$?
		\item Given $\theta \in \mathbb{R}$, show that
			\[
				\cos{\theta} = \frac{e^{i \theta} + e^{-i \theta}}{2} \quad \text{and} \quad \sin{\theta} = \frac{e^{i \theta} - e^{- i \theta}}{2 i}.
			\]
			These are also called Euler's identities.
		\item Use the complex exponential to derive trigonometric identities such as
			\[
				\cos{(\theta + \vartheta)} = \cos{\theta} \cos{\vartheta} - \sin{\theta} \sin{\vartheta},
			\]
			and then show that
			\begin{align*}
				2 \sin{\theta} \sin{\vartheta} & = \cos{(\theta - \vartheta)} - \cos{(\theta + \vartheta)} \\
				2 \sin{\theta} \cos{\vartheta} & = \sin{(\theta + \vartheta)} + \sin{(\vartheta - \vartheta)}.
			\end{align*}
	\end{enumerate}
\end{problem}

\begin{solution}
	\begin{enumerate}
		\item 
		\item Consider the $N$th partial sum of the left hand side ---
			\[
				 \sum_{n = 1}^{N} \frac{(z_{1} + z_{2})^{n}}{n!} = \sum_{n = 1}^{N} \frac{1}{n!} \sum_{k = 1}^{n} \binom{n}{k} z_{1}^{k} z_{2}^{n - k}
			\]
		\item Recall the power series
			\begin{align*}
				\cos{y} & = \sum_{n = 1}^{\infty} \frac{(-1)^{n}}{(2n)!} y^{2n} \\
				\sin{y}	& = \sum_{n = 1}^{\infty} \frac{(-1)^{n}}{(2n + 1)!} y^{2n + 1}.
			\end{align*}
			Computing the power series corresponding to $e^{iy}$ and picking out odd and even terms we see
			\begin{align*}
				e^{iy} & = \sum_{n = 1}^{\infty} \frac{(iy)^{n}}{n!} \\
				       & = \left ( 1 + i^{2} \frac{y^{2}}{2!} + i^{4} \frac{y^{4}}{4!} \dots \right ) + \left ( i y + i^{3} \frac{y^{3}}{3!} + \dots \right ) \\
				       & = \sum_{n = 1}^{\infty} \frac{(-1)^{n}}{(2n)!} y^{2n} + i \sum_{n = 1}^{\infty} \frac{(-1)^{n}}{(2n + 1)!} y^{2n + 1} \\
				       & = \cos{y} + i \sin{y}.
			\end{align*}
		\item By part (b), $e^{x + iy} = e^{x} e^{iy} = e^{x}(\cos{y} + i \sin{y})$. Thus,
			\[
				\lvert e^{x + i y} \rvert = \lvert e^{x} \rvert \lvert e^{iy} \rvert.
			\]
			Note that
			\begin{align*}
				\lvert e^{iy} \rvert^{2} & = \lvert \cos{y} + i \sin{y} \rvert^{2} \\
							 & = \cos^{2}{y} + \sin^{2}{y} \\
							 & = 1.
			\end{align*}
			Thus, the magnitude is just
			\[
				\lvert e^{x + iy} \rvert = \lvert e^{x} \rvert = e^{x}.
			\]
	\end{enumerate}
\end{solution}

\pagebreak

\begin{problem}[6.]
	Prove that if $f$ is a twice continuously differentiable function on $\mathbb{R}$ which is a solution of the equation
	\[
		f''(t) + c^{2} f(t) = 0
	\]
	then there exist constants $a$ and $b$ such that
	\[
		f(t) = a \cos{(c t)} + b \sin{(c t)}.
	\]
	This can be done by differentiating the two functions $g(t) = f(t) \cos{(c t)} - (1 / c) f'(t) \sin{(c t)}$ and $h(t) = f(t) \sin{(c t)} + (1 / c) f'(t) \cos{(ct)}$
\end{problem}

\begin{solution}
	We will first show that $g$ and $h$ are just constants and then show that $f(t) = g(t) \cos{(c t)} + h(t) \sin{(c t)}$.

	The first derivative of $g$ is
	\begin{align*}
		g'(t) & = f'(t) \cos{(c t)} - c f(t) \sin{(c t)} - (1/c) f''(t) \sin{(c t)} - (c / c) f'(t) \cos{(c t)} \\
		      & = f'(t) \cos{(c t)} - f'(t) \cos{(c t)} + (c^{2} / c) f(t) \sin{(c t)} - c f(t) \sin{(c t)} \\
		      & = 0.
	\end{align*}
	Thus, $g$ is constant. Note that in the second equality we use the fact that $f''(t) = - c^{2} f(t)$ from the differential equation.

	Similarly, the first derivative of $h$ is
	\begin{align*}
		h'(t) & = f'(t) \sin{(c t)} + c f(t) \cos{(c t)} + (1 / c) f''(t) \cos{(c t)} - (c / c) f'(t) \sin{(c t)} \\
		      & = f'(t) \sin{(c t)} - f'(t) \sin{(c t)} + c f(t) \cos{(c t)} - (c^{2} / c) f(t) \cos{(c t)} \\
		      & = 0.
	\end{align*}
	That is, $h$ is constant. Again, in the second equality we use the fact that $f''(t) = - c^{2} f(t)$.

	If we choose $a = g(0) = g(t)$ and $b = h(0) = h(t)$, we can check that
	\begin{align*}
		a \cos{(c t)} + b \sin{(c t)} & = g(t) \cos{(c t)} + h(t) \sin{(c t)} \\
					      & = f(t) \cos^{2}{(c t)} - (1 / c) f'(t) \sin{(c t)} \cos{(c t)} \\
					      & + f(t) \sin^{2}{(c t)} + (1 / c) f'(t) \cos{(c t)} \sin{(c t)} \\
					      & = f(t) (\cos^{2}{(c t)} + \sin^{2}{(c t)}) + f(t) (\cos{(c t)} \sin{(c t)} - \sin{(c t)} \cos{(c t)}) \\
					      & = f(t)
	\end{align*}
	as desired.
\end{solution} 

\pagebreak

\begin{problem}[8.]
	Suppose $F$ is a function on $(a, b)$ with two continuous derivatives. Show that whenever $x$ and $x + h$ belong to $(a, b)$, one may write
	\[
		F(x + h) = F(x) + h F'(x) + \frac{h^{2}}{2} F''(x) + h^{2} \varphi(h),
	\]
	where $\varphi(h) \to 0$ as $h \to 0$.

	Deduce that
	\[
		\frac{F(x + h) + F(x - h) - 2 F(x)}{h^{2}} \to F''(x)
	\]
	as $h \to 0$.
\end{problem}

\end{document}
